\chapter{Glossary}
\label{app:glossary}

Terms as this book uses them, in the order they are first introduced.

\begin{description}
\item[Driver] The process holding the \texttt{SparkSession}, running the
      application's own code, and coordinating everything else
      (\Cref{ch:shape}).
\item[Executor] A JVM process that runs tasks and holds data in memory or on
      disk for a running application (\Cref{ch:shape}).
\item[Job] The unit of work triggered by one action; made of one or more
      stages (\Cref{ch:shape}).
\item[Stage] A set of tasks that can run without a shuffle between them,
      bounded by wide-dependency shuffle boundaries (\Cref{ch:shape}).
\item[Task] The smallest unit of scheduled work: one stage's operators
      applied to one partition (\Cref{ch:shape}).
\item[Narrow dependency] A child partition depends on a known, fixed set of
      parent partitions --- no shuffle needed (\Cref{ch:shape}).
\item[Wide dependency] A child partition may depend on data from any parent
      partition --- requires a shuffle (\Cref{ch:shape}).
\item[Catalyst] Spark SQL's query optimizer: parse, analyze, optimize,
      plan, generate code (\Cref{ch:catalyst}).
\item[Whole-stage code generation] Fusing a chain of operators into one
      generated Java method via produce/consume
      (\Cref{sec:tungsten-wholestage}).
\item[Unified memory] The single JVM memory pool execution and storage
      share and borrow from, with a boundary either side may cross into the
      other's currently-idle space, but which only execution can enforce
      by evicting the borrower back out (\Cref{sec:mem-model}).
\item[Spill] Writing intermediate data (a sort buffer, a hash table) to disk
      because it does not fit in its memory budget (\Cref{ch:execmem}).
\item[Shuffle] Redistributing data across partitions between two stages, via
      a map-side write and a reduce-side fetch (\Cref{ch:shuffle}).
\item[Broadcast hash join] A join where one side is small enough to be sent
      whole to every executor, avoiding a shuffle of the large side
      (\Cref{ch:joins}).
\item[Sort-merge join] A join where both sides are shuffled and sorted by
      join key, then merged --- the robust default for two large sides
      (\Cref{ch:joins}).
\item[Adaptive Query Execution (AQE)] Re-planning after each shuffle using
      measured, rather than estimated, partition sizes
      (\Cref{ch:aqe}).
\item[Skew] One partition or key holding disproportionately more data than
      others, creating a straggler task (\Cref{sec:aqe-skew}).
\item[Partition pruning] Skipping whole partitions (or files) based on a
      filter resolved before reading data (\Cref{ch:pruning}).
\item[Dynamic partition pruning (DPP)] Partition pruning using a value
      discovered at runtime from the other side of a join
      (\Cref{sec:pruning-dpp}).
\item[Bucketing] Dividing data within a partition into a fixed number of
      files by hashing a column, to avoid a shuffle on a repeated join
      (\Cref{sec:pruning-bucketing}).
\item[Storage-partitioned join] A join with no shuffle because both sides
      are already co-located by the same bucket/partition transform (a
      partition transform --- see Hidden partitioning, below)
      (\Cref{sec:joins-spj}).
\item[Manifest] An immutable file listing a set of data or delete files with
      their partition values and column statistics (\Cref{ch:iceberg-spec}).
\item[Manifest list] A per-snapshot file listing the manifests that make up
      that snapshot, with per-manifest summary statistics
      (\Cref{ch:iceberg-spec}).
\item[Snapshot] One committed version of a table's state, with its own
      sequence number and manifest list (\Cref{sec:ice-snapshots}).
\item[Sequence number] A monotonically increasing number recording a
      snapshot's relative age, used to decide which delete files apply to
      which data files (\Cref{sec:ice-snapshots}).
\item[Position delete] A delete recorded as an exact (file, row-position)
      pair --- cheap to apply, requires the writer to know the row's
      position (\Cref{sec:ice-delete-files}).
\item[Equality delete] A delete recorded as a column-value match --- cheap
      to write, progressively more expensive to apply as they accumulate,
      because every read must check each one against every candidate row
      (\Cref{sec:ice-delete-files}).
\item[Copy-on-write (CoW)] A row-level update strategy that rewrites whole
      affected data files (\Cref{sec:ice-write-cow-mor}).
\item[Merge-on-read (MoR)] A row-level update strategy that records a delete
      file instead of rewriting data (\Cref{sec:ice-write-cow-mor}).
\item[Hidden partitioning] Iceberg partitioning by a transform of a source
      column, resolved from a plain filter with no partition column in the
      query text (\Cref{sec:ice-partitioning}).
\item[Catalog] The component mapping a table name to its current metadata
      file location, atomically swappable on commit
      (\Cref{sec:cat-what}).
\item[Optimistic concurrency] A commit protocol where writers do the
      expensive work assuming they will win, retrying only the cheap
      final compare-and-swap on conflict (\Cref{sec:cat-swap}).
\item[UnsafeRow] Tungsten's packed, off-heap-capable binary row format
      (\Cref{sec:tungsten-row}).
\item[DataSource V2 (DSv2)] The Spark API contract an external table
      source (Iceberg included) implements to plug into the planner
      (\Cref{ch:dsv2}).
\end{description}

\section*{Abbreviations}

\begin{table}[H]
\centering
\footnotesize
\renewcommand{\arraystretch}{0.9}
\begin{tabular}{@{}p{3.2cm}p{10.7cm}@{}}
\toprule
\textbf{Abbreviation} & \textbf{Full form} \\
\midrule
AQE & Adaptive Query Execution \\
CBO & Cost-Based Optimization \\
CoW / MoR & Copy-on-Write / Merge-on-Read \\
DAG & Directed Acyclic Graph \\
DPP & Dynamic Partition Pruning \\
DSv2 & DataSource V2 API \\
GC & Garbage Collection \\
GCS & Google Cloud Storage \\
GKE & Google Kubernetes Engine \\
HMS & Hive Metastore \\
JIT & Just-In-Time (compilation) \\
JVM & Java Virtual Machine \\
OLAP & Online Analytical Processing \\
OLTP & Online Transaction Processing \\
RDD & Resilient Distributed Dataset \\
RPC & Remote Procedure Call \\
SMJ & Sort-Merge Join \\
SQL & Structured Query Language \\
UDF & User-Defined Function \\
\bottomrule
\end{tabular}
\end{table}

\section*{Technologies}

\begin{table}[H]
\centering
\footnotesize
\renewcommand{\arraystretch}{0.9}
\begin{tabular}{@{}p{3.6cm}p{10.1cm}@{}}
\toprule
\textbf{Technology} & \textbf{What it is} \\
\midrule
Apache Spark & a distributed compute engine \\
Catalyst & Spark SQL's extensible query optimizer \\
Tungsten & Spark's binary in-memory execution back end \\
Apache Iceberg & an open table format for analytic datasets on object storage \\
Apache Parquet & a columnar file format \\
Hive Metastore & a catalog service mapping table names to metadata \\
Apache Kyuubi & a multi-tenant SQL gateway in front of Spark \\
Google Cloud Storage & an S3-compatible object storage system \\
GKE & Google's managed Kubernetes; the cluster manager here \\
async-profiler & a low-overhead JVM sampling profiler \\
\bottomrule
\end{tabular}
\end{table}
